\documentclass[11pt]{article}
\usepackage[margin=1in]{geometry}
\usepackage{amsmath,amssymb,amsthm}
\usepackage{mathtools}
\usepackage{hyperref}

\newtheorem{theorem}{Theorem}
\newtheorem{lemma}{Lemma}
\newtheorem{corollary}{Corollary}
\theoremstyle{definition}
\newtheorem{definition}{Definition}
\newtheorem{remark}{Remark}

\newcommand{\T}{\mathbb T}
\newcommand{\C}{\mathbb C}
\newcommand{\Lp}{L^p}
\newcommand{\norm}[1]{\left\lVert #1\right\rVert}
\newcommand{\supp}{\operatorname{supp}}

\title{Weak and Strong \(T_{L^p}\) for Second Countable Locally Compact Groups}
\author{run \texttt{fa\_banach\_001}, agent \texttt{agent\_super\_00}}
\date{2026-06-30}

\begin{document}
\maketitle

\section*{Status}

This is a likely valid strong partial result for arXiv:2310.18136.  The
source paper records that weak property \(T_{L^p}\) and property \(T_{L^p}\)
are known to coincide for countable discrete groups, and says that the
corresponding locally compact question is open.  The result below proves the
second countable locally compact case.  Arbitrary non-second-countable locally
compact groups are not claimed.

\section*{Source Question}

The source is E. M. Elki\ae r and S. Pooya, \emph{Property (T) for Banach
algebras}, arXiv:2310.18136 \cite{ElkiaerPooya}.  In the remark following
Theorem 4.3, after discussing \(F^*_{\lambda_p}(G)\), it states that for
countable discrete groups the equivalence between weak property
\(T_{L^p}\) and property \(T_{L^p}\) is known, but that it is open whether
this is true for general locally compact groups.

\section*{Theorem}

\begin{theorem}\label{thm:main}
Let \(G\) be a second countable locally compact group and let
\(1\le p<\infty\).  Then \(G\) has weak property \(T_{L^p}\) if and only if
it has property \(T_{L^p}\).
\end{theorem}

Here \(L^p\) denotes the class of complex \(L^p\)-spaces over sigma-finite
measure spaces, and representations are strongly continuous isometric
representations.

\section*{Proof Intuition}

The discrete proof of Elki\ae r \cite{ElkiaerWeak} uses circle-valued
measurable cocycles.  If coboundaries were closed for every ergodic p.m.p.
action, the open mapping theorem for Polish groups would force strong
ergodicity.  For a group without property \(T\), Connes-Weiss gives an
ergodic action that is not strongly ergodic, hence a cocycle in the closure of
the coboundaries but not itself a coboundary.  This cocycle twists the
Koopman representation on \(L^p\); approximating coboundaries give almost
invariant unit vectors, while non-coboundary prevents invariant vectors.

For locally compact second countable \(G\), the same mechanism works after
replacing the countable product topology on cocycles by compact-open
convergence in measure.  Compact sets play exactly the role of finite subsets
in the definition of almost invariant vectors.

\section*{Cocycles}

Let \(G\curvearrowright (\Omega,\mu)\) be an ergodic probability-measure
preserving action on a standard probability space.  We write
\[
   (g f)(\omega)=f(g^{-1}\omega)
\]
for the induced action on measurable functions.  Let \(L^0(\Omega;\T)\) be
equipped with convergence in measure; equivalently, for circle-valued
functions, with any \(L^p\)-metric.

A \(\T\)-valued \(1\)-cocycle is a continuous map
\[
   c:G\to L^0(\Omega;\T),\qquad g\mapsto c_g,
\]
where continuity is for convergence in measure, satisfying
\[
   c_{gh}(\omega)=c_g(\omega)c_h(g^{-1}\omega)
\]
for all \(g,h\in G\), almost everywhere.  Let \(Z^1(G\curvearrowright
\Omega;\T)\) be the group of these cocycles with the compact-open topology:
\[
   c_i\to c
   \quad\Longleftrightarrow\quad
   \sup_{g\in K}\norm{c_{i,g}-c_g}_{L^p(\mu)}\to0
   \quad\hbox{for every compact }K\subset G.
\]
Because \(G\) is second countable locally compact and \(L^0(\Omega;\T)\) is
Polish, this compact-open topology makes \(Z^1\) a Polish group; the cocycle
identity defines a closed subgroup of the Polish group \(C(G,L^0(\Omega;\T))\).

For \(\phi\in L^0(\Omega;\T)\), define
\[
   b_\phi(g,\omega)=\phi(\omega)\overline{\phi(g^{-1}\omega)} .
\]
The set \(B^1\) of such coboundaries is a subgroup of \(Z^1\).  The map
\(\beta:L^0(\Omega;\T)\to B^1\), \(\beta(\phi)=b_\phi\), is a continuous
surjective homomorphism.  Continuity into the compact-open topology follows
from strong continuity of the action on \(L^p(\Omega)\), uniformly on compact
subsets, and from the fact that all functions involved are \(\T\)-valued.
Since the action is ergodic, the kernel consists precisely of the constants.

\begin{lemma}\label{lem:closed-strong}
If \(B^1(G\curvearrowright\Omega;\T)\) is closed in \(Z^1(G\curvearrowright
\Omega;\T)\), then the action is strongly ergodic.
\end{lemma}

\begin{proof}
If \(B^1\) is closed, then it is Polish.  The open mapping theorem for Polish
groups applied to \(\beta:L^0(\Omega;\T)\to B^1\) shows that \(\beta\) is
open.  Fix \(0<\varepsilon<1/2\) and set
\[
   U_\varepsilon=\{\phi\in L^0(\Omega;\T):
       \inf_{z\in\T}\norm{\phi-z}_{L^p(\mu)}<\varepsilon\}.
\]
Then \(\beta(U_\varepsilon)\) is an open neighborhood of the identity in
\(B^1\).  Hence there are a compact set \(K\subset G\) and \(\delta>0\) such
that, whenever \(\phi\in L^0(\Omega;\T)\) satisfies
\[
   \sup_{g\in K}\norm{b_\phi(g)-1}_{L^p(\mu)}<\delta,
\]
we have \(b_\phi\in \beta(U_\varepsilon)\).  This means that after multiplying
\(\phi\) by a constant in \(\T\), it is within \(\varepsilon\) of \(1\) in
\(L^p\).

Suppose the action were not strongly ergodic.  Then there is a net of
measurable sets \(A_i\) with \(\mu(A_i)=1/2\) such that
\[
   \sup_{g\in K}\mu(gA_i\triangle A_i)\to0
\]
for every compact \(K\subset G\).  Put
\(\phi_i=1_{A_i}-1_{\Omega\setminus A_i}\).  Then
\[
   \norm{b_{\phi_i}(g)-1}_{L^p}^p
   =2^p\mu(gA_i\triangle A_i),
\]
so \(b_{\phi_i}\to1\) in \(B^1\).  By the previous paragraph, \(\phi_i\) is
eventually within \(\varepsilon\) of a constant circle-valued function.  This
is impossible for \(\varepsilon<1\), since every \(\phi_i\) takes the two
values \(1\) and \(-1\) on sets of measure \(1/2\).  Thus the action is
strongly ergodic.
\end{proof}

\begin{lemma}\label{lem:twist}
Let \(c\in Z^1(G\curvearrowright\Omega;\T)\).  Define
\[
   (\pi_c(g)f)(\omega)=c_g(\omega)f(g^{-1}\omega),\qquad f\in L^p(\Omega).
\]
Then \(\pi_c\) is a strongly continuous isometric representation of \(G\) on
\(L^p(\Omega)\).

Moreover:
\begin{enumerate}
\item if \(c\in\overline{B^1}\), then \(\pi_c\) has almost invariant unit
vectors;
\item if \(c\notin B^1\), then \(\pi_c\) has no non-zero invariant vector.
\end{enumerate}
\end{lemma}

\begin{proof}
The cocycle identity gives the representation identity, and the values in
\(\T\) make the operators isometric.  Strong continuity follows from
continuity of \(g\mapsto c_g\) in measure and continuity of the Koopman
representation; bounded convergence upgrades this to \(L^p\)-continuity on
circle-valued multipliers.

Assume \(c\in\overline{B^1}\), and choose a net \(\phi_i\in L^0(\Omega;\T)\)
with \(b_{\phi_i}\to c\) in compact-open topology.  Each \(\phi_i\) is a unit
vector in \(L^p(\Omega)\), and for every compact \(K\subset G\),
\[
\begin{aligned}
   \sup_{g\in K}\norm{\pi_c(g)\phi_i-\phi_i}_{L^p}
   &=\sup_{g\in K}\norm{
       c_g(\omega)\phi_i(g^{-1}\omega)-\phi_i(\omega)}_{L^p_\omega} \\
   &=\sup_{g\in K}\norm{c_g-b_{\phi_i}(g)}_{L^p}\to0.
\end{aligned}
\]
Thus \(\pi_c\) has almost invariant unit vectors.

Conversely, suppose \(0\ne f\in L^p(\Omega)\) is invariant.  Then
\(|f|\) is invariant under the original p.m.p. action.  Ergodicity implies
that \(|f|\) is almost everywhere constant and non-zero.  Hence
\(\phi=f/|f|\) is circle-valued and
\[
   c_g(\omega)\phi(g^{-1}\omega)=\phi(\omega)
\]
for every \(g\), almost everywhere.  Equivalently \(c=b_\phi\), contrary to
\(c\notin B^1\).  Therefore no non-zero invariant vector exists when \(c\) is
not a coboundary.
\end{proof}

\begin{proof}[Proof of Theorem \ref{thm:main}]
Property \(T_{L^p}\) plainly implies weak property \(T_{L^p}\).  For the
converse, assume that \(G\) does not have property \(T_{L^p}\).  By the
Bader-Furman-Gelander-Monod theorem for second countable locally compact
groups \cite{BaderFurmanGelanderMonod}, this means that \(G\) does not have
Kazhdan's property \(T\).

By the locally compact Connes-Weiss-Schmidt strong-ergodicity
characterization of property \(T\), there is an ergodic p.m.p. action of
\(G\) on a standard probability space which is not strongly ergodic.  By
Lemma \ref{lem:closed-strong}, \(B^1\) is not closed in \(Z^1\).  Choose
\[
   c\in \overline{B^1}\setminus B^1 .
\]
Lemma \ref{lem:twist} gives a strongly continuous isometric representation
\(\pi_c\) on the \(L^p\)-space \(L^p(\Omega)\) with almost invariant unit
vectors and no non-zero invariant vector.  Hence \(G\) fails weak property
\(T_{L^p}\).  This proves that weak property \(T_{L^p}\) implies property
\(T_{L^p}\).
\end{proof}

\section*{Consequence for the Source Paper}

The source paper proves that for discrete groups \(F^*_{\lambda_p}(\Gamma)\)
detects \(T_{L^p}\).  It also remarks that, for non-discrete IN groups and
\(p\ne2\), the same method gives only weak \(T_{L^p}\) for the group.  Theorem
\ref{thm:main} upgrades that weak conclusion to \(T_{L^p}\) whenever the IN
group is second countable.

\section*{Limitations and Novelty Check}

This packet does not treat arbitrary non-second-countable locally compact
groups.  The proof uses second countability to make the compact-open cocycle
space Polish and to invoke the usual standard probability-space form of the
strong-ergodicity characterization.

Cheap run indexes had no prior packet for arXiv:2310.18136.  The source paper
cites arXiv:2403.05312 \cite{ElkiaerWeak} for countable discrete groups and
explicitly records the locally compact case as open.  Bounded web searches on
2026-06-30 for exact phrases including
\[
\text{``weak property'' ``}T_{L^p}\text{'' ``locally compact''}
\]
and
\[
\text{``}F^*_{\lambda_p}(G)\text{'' ``property'' ``T'' ``non-discrete''}
\]
found no later settlement.

\section*{Human Review Recommendation}

Review as a likely valid strong partial result.  The main points to check are
the compact-open Polish topology on \(Z^1\), the uniform-on-compact continuity
of the coboundary map, and the exact locally compact form of the
Connes-Weiss-Schmidt theorem used in the final proof.

\begin{thebibliography}{9}

\bibitem{BaderFurmanGelanderMonod}
U. Bader, A. Furman, T. Gelander, and N. Monod,
\emph{Property (T) and rigidity for actions on Banach spaces},
Acta Math. 198 (2007), 57--105.

\bibitem{ConnesWeiss}
A. Connes and B. Weiss,
\emph{Property T and asymptotically invariant sequences},
Israel J. Math. 37 (1980), 209--210.

\bibitem{ElkiaerWeak}
E. M. Elki\ae r,
\emph{Weak property \(T_{L^p}\) for discrete groups},
arXiv:2403.05312.

\bibitem{ElkiaerPooya}
E. M. Elki\ae r and S. Pooya,
\emph{Property (T) for Banach algebras},
arXiv:2310.18136.

\bibitem{Schmidt}
K. Schmidt,
\emph{Asymptotically invariant sequences and an action of \(SL(2,\mathbb Z)\)
on the 2-sphere},
Israel J. Math. 37 (1980), 193--208.

\end{thebibliography}

\end{document}
